\section{Evaluation Methodology}\label{sec:methodology}
% Thesis source: sections 6.1--6.6.
\begin{itemize}[leftmargin=*,itemsep=0.35em]
  \item Describe the repeated scenario-based comparison across model and orchestration configurations, followed by detailed tests of a selected configuration.
  \item Define functional success, correct tool behaviour, result correctness, response time, token usage and model-call cost. Explain aggregation and variability without disclosing internal costs or confidential measurements.
  \item Use the existing rule-based REST integration only as a qualitative reference; it does not implement an identical conversational task.
  \item If space permits, briefly describe exploratory user feedback, normal and exceptional cases, and robustness or concurrency tests. Keep their methods abstract and omit raw logs and full questionnaires.
\end{itemize}

\section{Results and Discussion}\label{Sec:Evaluation}
% Thesis source: chapters 6--7. No internal results are reproduced in this outline.
\begin{itemize}[leftmargin=*,itemsep=0.35em]
  \item Reserve space for a compact comparison of reliability, latency and resource use. Include only results cleared for external disclosure; check each later figure, table and result for confidential information.
  \item Discuss whether the separate planner adds value and the trade-off between fine-grained tools and bundled backend operations.
  \item Relate the findings to the research question and distinguish observations for the investigated prototype from general recommendations.
  \item Include exploratory feedback or exception-case findings only if space permits and their disclosure is permitted; explain their evidential limits.
\end{itemize}

\subsection*{Template examples: figures, tables and acronyms}
\textit{The following examples illustrate formatting and do not report results of the scheduling study.}

We evaluate the performance of our methods on the meshes of Fig.~\ref{fig:meshes}.
    
\begin{figure*}
    \centering%
    \def\svgwidth{\textwidth}%
    \fontsize{6pt}{5pt}\selectfont%
    \includesvg{figs/meshes}%
    \caption[Test Meshes]{
    $V$ denotes the number of vertices of the input mesh, $T$ the number of triangles and $M$ the number of meshlets coming from Meshoptimizer.
    Meshlets are visualized with randomized colors.
    }\label{fig:meshes}%
\end{figure*}%

\begin{table}
    \caption
    [Comparison for the \textit{Rock} mesh.]
    {    
    We compare the optimal Gurobi and SCIP solutions against the sub-optimal stuff.
    The CPU computation uses one thread per meshlet and was measured on an AMD Ryzen 9 7950X (16C/32T).    
    }\label{tab:StripifyTable}    
    {
    \centering    
    \footnotesize
\begin{tblr}[b]{lccc}
    \toprule                                        
                     & \textbf{Gurobi} & \textbf{SCIP}       & ETA\\
\midrule                                    
\textbf{computation time}     & 122.46\,s    & 8,122.56\,s & 10.52\,s   \\
{GTS render time} & 1.58\,ms     & 1.58\,ms    & 0.59\,ms \\
{GTS restarts}    & 123,784      & 123,784     & 185,480\\
\textbf{degenerate triangles} & 2,132      & 495,136     & 741,920\\
\textbf{additional meshlets}  & 1 / 12412   & 123  & 4 / 75,413\\
\bottomrule                                    
\end{tblr}


    }
\end{table}


\begin{figure}
    \includesvg[inkscapelatex=true]{figs/PDFTexFigure}
    \caption[Vector graphics with fonts rendered by latex]{You can use Latex to render the fonts.}\label{fig:latexfonts}
\end{figure}

\begin{figure}
    \includesvg[inkscapelatex=false]{figs/SVGFigure}
    \caption[Vector graphics with fonts directly taken from the file]{You can also use the fonts produced by the SVG file.}\label{fig:svgfonts}
\end{figure}

\begin{figure}
    \centering
    \includegraphics[width=\columnwidth]{imgs/img.png}
    \caption[Pixel graphics]{You can also directly include pixel graphics.}\label{fig:pixel}
\end{figure}

Tab.~\ref{tab:StripifyTable} and Tab.~\ref{tab:ListOfSymbols} compare our optimal 
method achieved with our other method
of Sec.~\ref{Sec:MainPart} with Gurobi and SCIP against the sub-optimal strip of our 
other implementation.
See Figs.~\ref{fig:latexfonts}, Figs.~\ref{fig:svgfonts}, and Figs.~\ref{fig:pixel} for different types of figures.
Here we test the acronyms.

% Test acronyms
First time we use a \ac{GPU}.
Next time, we use a \ac{GPU}.
Now, we even have multiple \acp{GPU}.
